%==============================================================================
%  MATH 231 -- Linear Algebra I (Queens College, Fall 2026)
%  COMPUTATIONAL UNIT, PART 2 -- k-Nearest Neighbours
%
%  Section numbering runs CONTINUOUSLY across the parts of the computational
%  unit, matching the vector unit's scheme:
%
%      Part 1  MATH231_computational_1_operation_counting        Sections 1-4
%      Part 2  MATH231_computational_2_k_nearest_neighbours      Sections 5-10
%      Part 3  MATH231_computational_3_k_means_clustering        Sections 11-15
%
%  Hence \setcounter{section}{4}: this part opens at Section 5.  References to
%  Sections 1-4 point into Part 1.  Cross-references OUT of the unit still name
%  the unit ("the vector unit, Sec. 14.2"), because the vector unit also has a
%  Sec. 14 and the matrix unit also has a Sec. 1.
%
%  STATUS: complete.  Sections 5-10.
%
%  AUDIENCE: distributed to students.  Keep it free of instructor-facing
%  apparatus -- time budgets, pacing notes, teaching cues.
%
%  THE CLUSTERING QUESTION, Sec. 9.  "k-NN clustering" is a phrase in wide
%  circulation and it names two different things.  k-NN itself is SUPERVISED
%  classification -- it needs labels and it does not cluster.  What legitimately
%  bears the name is the k-NN GRAPH, whose connected components (or spectral
%  cuts) are used as clusters, and which is the substrate of DBSCAN, spectral
%  clustering and UMAP.  Sec. 9 gives both: a flagged warning drawing the
%  distinction, then a worked 1-NN graph on six points whose two components are
%  the two clusters.  Do not quietly "fix" this by calling k-NN a clustering
%  algorithm, and do not drop Sec. 9 -- the phrase will be met in the wild.
%
%  COUNTS REUSED FROM PART 1.  Sec. 7 does not re-derive the cost of a norm; it
%  cites Sec. 1.4 (2n operations) and adds the n subtractions that form the
%  difference vector, giving 3n for a Euclidean distance.  If Sec. 1.4's count
%  ever changes, Sec. 7.1 changes with it.  The squared-distance shortcut of
%  Sec. 7.5 is deliberately framed as an instance of Sec. 1.6's "do not take a
%  square root you are going to square".
%
%  THE WORKED EXAMPLE in Sec. 6.5 is built so that k=1 and k=3 DISAGREE: the
%  point a4 = (4,5) is an A-labelled outlier sitting inside the B region, and it
%  is the single nearest neighbour of the query.  All eight distances are
%  distinct, deliberately, so no tie-breaking rule is needed there; ties get
%  their own treatment in Sec. 8.4.  Changing any coordinate risks creating a
%  tie and destroying the k=1 vs k=3 contrast, which Sec. 8.1 then builds on.
%
%  Compile with pdflatex (standard TeX Live packages only; uses tikz).
%==============================================================================
\documentclass[11pt]{article}

\usepackage[margin=1in]{geometry}
\usepackage{amsmath,amssymb,amsthm}
\usepackage[dvipsnames]{xcolor}
\usepackage{enumitem}
\usepackage{booktabs}
\usepackage{array}
\usepackage[hidelinks]{hyperref}
\usepackage{titlesec}
\usepackage{needspace}
\usepackage{tikz}
\usetikzlibrary{arrows.meta,shapes.geometric}

%------------------------------------------------------------------ style ----
\titleformat{\section}{\large\bfseries\sffamily}{\thesection.}{0.6em}{}
\titleformat{\subsection}{\normalsize\bfseries\sffamily}{\thesubsection.}{0.5em}{}
\titlespacing*{\section}{0pt}{14pt}{6pt}

\theoremstyle{plain}
\newtheorem{thm}{Theorem}[section]

\theoremstyle{definition}
\newtheorem{defn}{Definition}[section]
\newtheorem{ex}{Example}[section]
\theoremstyle{remark}
\newtheorem*{rmk}{Remark}
\newtheorem*{warn}{A common confusion}

\newcommand{\lookahead}{\par\smallskip\noindent
  \textbf{\sffamily\textcolor{RoyalBlue}{Looking ahead.}}\ \ignorespaces}
\newcommand{\lookback}{\par\smallskip\noindent
  \textbf{\sffamily\textcolor{ForestGreen}{From the vector unit.}}\ \ignorespaces}

%--------------------------------------------------------------- notation ----
\newcommand{\R}{\mathbb{R}}
\newcommand{\F}{\mathbb{F}}
\newcommand{\vc}[1]{\mathbf{#1}}
\newcommand{\norm}[1]{\lVert #1\rVert}
\newcommand{\D}{\mathcal{D}}
\newcommand{\Nk}{\mathcal{N}_k}
\DeclareMathOperator*{\argmin}{arg\,min}
\DeclareMathOperator*{\argmax}{arg\,max}

%--------------------------------------------------------------- pseudocode --
\newcommand{\kw}[1]{\textsf{\bfseries #1}}
\newcommand{\ind}{\hspace*{1.3em}}
\newcommand{\indd}{\hspace*{2.6em}}
\newenvironment{algo}[2]{%
  \par\medskip\needspace{8\baselineskip}%
  \begin{center}\begin{minipage}{0.94\textwidth}%
  \hrule height 0.9pt \vspace{3pt}%
  \noindent{\sffamily\bfseries #1}\par\vspace{2pt}%
  \noindent{\small #2}\par\vspace{3pt}%
  \hrule height 0.4pt \vspace{5pt}%
  \begin{list}{\textbf{\arabic{algoline}:}}%
    {\usecounter{algoline}\setlength{\leftmargin}{2.3em}%
     \setlength{\labelwidth}{1.8em}\setlength{\labelsep}{0.5em}%
     \setlength{\itemsep}{2pt}\setlength{\topsep}{0pt}%
     \setlength{\parsep}{0pt}}%
}{%
  \end{list}\vspace{5pt}\hrule height 0.9pt%
  \end{minipage}\end{center}\medskip%
}
\newcounter{algoline}

%------------------------------------------------------------- tikz helpers --
\tikzset{
  opbox/.style    = {draw=gray!65, rounded corners=2pt, fill=gray!8,
                     inner sep=4pt, font=\small, minimum height=7.5mm},
  opar/.style     = {-{Stealth[length=2.4mm,width=1.8mm]}, gray!80, thick},
  axisline/.style = {->, gray!75, thin},
  ptA/.style      = {circle, fill=RoyalBlue!75, draw=RoyalBlue!90,
                     inner sep=0pt, minimum size=5.4pt},
  ptB/.style      = {rectangle, fill=BrickRed!75, draw=BrickRed!90,
                     inner sep=0pt, minimum size=5.2pt},
  ptQ/.style      = {star, star points=5, fill=ForestGreen!85,
                     draw=ForestGreen, inner sep=0pt, minimum size=8pt}
}

\title{\vspace{-1.2cm}\textbf{\sffamily MATH 231 -- Linear Algebra I}\\[4pt]
       {\large\sffamily The Computational Unit, Part 2 \textbullet\
        $k$-Nearest Neighbours}\\[2pt]
       {\normalsize\sffamily Kiely Hall 326 \textbullet\
        Instructor: Kevin Bedoya}}
\author{}
\date{}

\begin{document}
\setcounter{section}{4}   % this part opens at Section 5; see the header note
\maketitle
\vspace{-1.6cm}

\noindent\textbf{About this document.} This is Part 2 of the computational
unit, and the first of two on algorithms that do nothing but measure distances
between vectors. Everything it needs was built in the vector unit; what is new
is putting those tools to work on a problem, and then asking what the work
costs.

\begin{center}
\renewcommand{\arraystretch}{1.2}
\begin{tabular}{@{}l l l@{}}
\toprule
\textbf{Part} & \textbf{Sections} & \textbf{Subject} \\
\midrule
1 & \S1--\S4    & operation counting, Big-O, floating point, error \\
2 & \S5--\S10   & $k$-nearest neighbours \\
3 & \S11--\S15  & $k$-means clustering \\
\bottomrule
\end{tabular}
\end{center}

\begin{warn}[Whose \S5 is this?]
Section numbering runs continuously through the computational unit, so a bare
``\S1.4'' below means \S1.4 \emph{of this unit} --- Part 1. The vector unit and
the matrix unit have their own numbering, and references into them always say
so, as in ``the vector unit, \S14.2''.
\end{warn}

\medskip
\noindent\textbf{Assumed.} From the vector unit: $\R^{n}$ and its arithmetic
(\S24); the Euclidean norm (\S2) and distance $d(\vc{x},\vc{y})=\norm{\vc{x}-\vc{y}}$
(\S14.2); the metric axioms \textbf{M1}--\textbf{M4} (\S14.1); the family of
$p$-norms and the Manhattan and Chebyshev metrics (\S14.3--\S14.6); cosine
similarity (\S11.4). From Part 1 of this unit: what an operation is (\S1.2), the
cost of a norm (\S1.4), Big-O (\S2.7), and FLOPs (\S4.6).

\medskip
\noindent\textbf{What you should be able to do afterwards.} State the
classification problem in vector terms; describe the $k$-nearest neighbours rule
and run it by hand; write its pseudocode; count its operations and give its
complexity; explain what changes as $k$ grows and why feature scaling matters;
distinguish $k$-NN from clustering and build a $k$-NN graph; and name the
standard ways the algorithm is made fast in practice.

%==============================================================================
\section{Data as vectors}
%==============================================================================
%------------------------------------------------------------------------------
\subsection{Turning a thing into a point}
%------------------------------------------------------------------------------
Everything in this part rests on one move, and the move is so routine in
practice that it is easy to miss how much it assumes. Take some collection of
objects --- patients, photographs, emails, houses for sale --- and describe each
one by a fixed list of $n$ numbers. That list is a vector in $\R^{n}$, and the
object has become a point.

\begin{defn}[Feature vector]
A \textbf{feature vector} is a vector $\vc{x}=[\,x_1,\dots,x_n\,]\in\R^{n}$ whose
components are $n$ measured attributes of a single object, recorded in a fixed
order. Each component is a \textbf{feature}, and $n$ is the
\textbf{dimension} of the data.
\end{defn}

\noindent A house might become $[\,\text{area},\ \text{bedrooms},\
\text{age},\ \text{distance to station}\,]\in\R^{4}$. A grayscale image
$28\times28$ pixels becomes the list of its $784$ brightness values, a point in
$\R^{784}$. A sentence, run through a modern language model, becomes a vector of
a few hundred or a few thousand numbers.

\begin{defn}[Dataset]
A \textbf{labelled dataset} is a finite collection
\[
  \D = \bigl\{(\vc{x}_1,y_1),\ (\vc{x}_2,y_2),\ \dots,\ (\vc{x}_N,y_N)\bigr\},
  \qquad \vc{x}_i\in\R^{n},
\]
of $N$ feature vectors, each carrying a \textbf{label} $y_i$. When the labels
come from a finite set of categories the problem is \textbf{classification};
when they are real numbers it is \textbf{regression}.
\end{defn}

\noindent Two numbers will recur throughout, and they measure different things:
$N$ is how \emph{many} points there are, and $n$ is how \emph{long} each one is.
Confusing them is the most common error in a first complexity analysis.

\begin{rmk}[What the move assumes]
Writing an object as a point in $\R^{n}$ commits to something: that the
arithmetic of $\R^{n}$ means something about the objects. Specifically, it
commits to \emph{distance being meaningful} --- to the claim that two houses
whose feature vectors are close really are similar houses. That is a modelling
assumption, not a mathematical fact, and \S8.2 is about a way it commonly fails.
\end{rmk}

%------------------------------------------------------------------------------
\subsection{The classification problem}
%------------------------------------------------------------------------------
\begin{quote}
Given a labelled dataset $\D$ and a new, unlabelled point $\vc{q}\in\R^{n}$ ---
the \textbf{query} --- predict its label.
\end{quote}

\noindent The dataset is what we have been told; the query is what we are asked
about. A \textbf{classifier} is any rule that turns the first into an answer
about the second.

%------------------------------------------------------------------------------
\subsection{The idea in one sentence}
%------------------------------------------------------------------------------
Here is the whole of $k$-nearest neighbours, before any notation:

\begin{center}
\emph{To label a new point, look at the labelled points nearest to it, and let
them vote.}
\end{center}

\noindent That is an idea, not yet an algorithm, and turning it into one
requires answering exactly three questions.

\begin{enumerate}[label=\textbf{Q\arabic*.},leftmargin=2.6em,itemsep=4pt]
  \item What does \emph{nearest} mean? --- \S5.4.
  \item How many neighbours vote? --- that is the $k$, \S6.2.
  \item How do they vote? --- \S6.3.
\end{enumerate}

%------------------------------------------------------------------------------
\subsection{What ``nearest'' means}
%------------------------------------------------------------------------------
\lookback The vector unit answered Q1 in advance. A \emph{metric} is a function
$d(\vc{x},\vc{y})$ measuring distance, subject to the four axioms
\textbf{M1}--\textbf{M4} of \S14.1, and the unit built a whole family of them.
The default is Euclidean distance,
\[
  d_2(\vc{x},\vc{y}) = \norm{\vc{x}-\vc{y}}_2
  = \sqrt{\sum_{i=1}^{n}(x_i-y_i)^{2}} \qquad (\S14.2),
\]
but nothing in the algorithm requires it. The Manhattan distance $d_1$
(\S14.3), the Chebyshev distance $d_\infty$ (\S14.4), and any $p$-norm distance
(\S14.6) all satisfy the axioms and all may be substituted directly.

Which one to use is a modelling decision. Manhattan distance is natural when the
features are independent counts and you want each to contribute in proportion to
its own deviation; Euclidean is natural when the features are spatial
coordinates. For text and for high-dimensional embeddings, practitioners very
often use \emph{cosine} distance instead,
\[
  d_{\cos}(\vc{x},\vc{y}) = 1 - \cos(\vc{x},\vc{y})
  = 1 - \frac{\langle\vc{x},\vc{y}\rangle}{\norm{\vc{x}}\norm{\vc{y}}}
  \qquad (\text{the vector unit, }\S11.4),
\]
which compares \emph{direction} and ignores magnitude entirely --- appropriate
when a long document and a short one on the same topic should count as similar.

\begin{rmk}
$d_{\cos}$ is not a metric: it fails the triangle inequality \textbf{M4}. This
is worth knowing and, for $k$-NN, worth not worrying about --- the algorithm
only ever \emph{sorts} by the distance function, and sorting needs an ordering,
not an axiom. Some of the fast methods in \S10.2 do rely on the triangle
inequality, and those are the ones that stop working.
\end{rmk}

%==============================================================================
\section{The algorithm}
%==============================================================================
%------------------------------------------------------------------------------
\subsection{Walking through it}
%------------------------------------------------------------------------------
Fix a metric $d$ and an integer $k\ge1$. Given the query $\vc{q}$:

\begin{enumerate}[label=\textbf{Step \arabic*.},leftmargin=4.2em,itemsep=5pt]
  \item \textbf{Measure.} Compute the distance from $\vc{q}$ to every one of the
        $N$ points in $\D$. This produces $N$ numbers
        $d(\vc{q},\vc{x}_1),\dots,d(\vc{q},\vc{x}_N)$ and nothing else.
  \item \textbf{Select.} Find the $k$ smallest of those numbers, and remember
        which points they belonged to. Call that set of points the
        \textbf{$k$-neighbourhood} of $\vc{q}$, written $\Nk(\vc{q})$.
  \item \textbf{Vote.} Look at the labels of those $k$ points. Predict whichever
        label occurs most often.
\end{enumerate}

\noindent That is the entire algorithm. Notice what is \emph{not} in it: there
is no model, no equation fitted to the data, no parameters estimated. The
dataset is the model. Everything happens when a query arrives.

\begin{center}
\begin{tikzpicture}[scale=0.74]
  \draw[axisline] (-0.5,0) -- (9.8,0) node[right,font=\scriptsize,black]{$x_1$};
  \draw[axisline] (0,-0.5) -- (0,9.9) node[above,font=\scriptsize,black]{$x_2$};
  % the two neighbourhood balls, drawn just outside the 1st and 3rd neighbour
  % (which sit at distance 1 and 2.236) so that they visibly enclose them
  \draw[ForestGreen!70,dashed] (5,5) circle (1.15);
  \draw[ForestGreen!45,dashed] (5,5) circle (2.40);
  % class A points
  \foreach \x/\y in {1/1, 2/3, 2/1, 4/5} { \node[ptA] at (\x,\y) {}; }
  % class B points
  \foreach \x/\y in {6/4, 7/6, 6/8, 8/5} { \node[ptB] at (\x,\y) {}; }
  \node[ptQ] at (5,5) {};
  \node[font=\scriptsize,ForestGreen!85] at (3.72,4.05) {$k=1$};
  \node[font=\scriptsize,ForestGreen!65] at (3.00,3.15) {$k=3$};
  \node[above,font=\scriptsize,ForestGreen!85] at (5.0,5.28) {$\vc{q}$};
  \node[above left,font=\scriptsize,RoyalBlue] at (4.08,5.05) {$\vc{a}_4$};
  \node[below right,font=\scriptsize,BrickRed] at (6.05,3.95) {$\vc{b}_1$};
  % legend
  \node[ptA] at (0.7,9.4) {}; \node[right,font=\scriptsize] at (0.95,9.4) {class $A$};
  \node[ptB] at (3.4,9.4) {}; \node[right,font=\scriptsize] at (3.65,9.4) {class $B$};
  \node[ptQ] at (6.1,9.4) {}; \node[right,font=\scriptsize] at (6.4,9.4) {query};
\end{tikzpicture}
\end{center}

\noindent The two dashed circles are the picture of Step~2: growing a ball
around $\vc{q}$ until it contains exactly $k$ points. The inner circle catches
one point, which is blue; the outer catches three, of which two are red. So the
two settings of $k$ disagree, and \S6.5 works out why in numbers.

%------------------------------------------------------------------------------
\subsection{The $k$}
%------------------------------------------------------------------------------
\begin{defn}[$k$-neighbourhood]
The \textbf{$k$-neighbourhood} of a query $\vc{q}$, written $\Nk(\vc{q})$, is
the set of the $k$ points of the dataset $\D$ lying closest to $\vc{q}$ under
the metric $d$.
\end{defn}

\noindent Two extreme settings frame the choice.

\begin{itemize}[leftmargin=1.4em,itemsep=4pt]
  \item $k=1$ is the \textbf{nearest neighbour} rule: copy the label of the
        single closest point. It follows the data exactly --- and therefore
        follows its mistakes exactly too. One mislabelled or unusual training
        point creates a small region around itself where every query gets the
        wrong answer.
  \item $k=N$ makes every query return the most common label in the entire
        dataset, regardless of where the query is. The distances have stopped
        mattering.
\end{itemize}

\noindent Useful values live between, and \S8.1 says how to think about
choosing. For a two-class problem it is customary to take $k$ \emph{odd}, which
makes a tied vote impossible.

%------------------------------------------------------------------------------
\subsection{Pseudocode}
%------------------------------------------------------------------------------
Write $\mathbf{1}[\,P\,]$ for the \emph{indicator} of a statement $P$: it is $1$
when $P$ is true and $0$ when $P$ is false. Then counting how many neighbours
carry the label $c$ is a sum of indicators, and the vote is an $\argmax$ over
the possible labels.

\begin{algo}{$k$-NN classification}
{\emph{Input:} a labelled dataset $\D=\{(\vc{x}_i,y_i)\}_{i=1}^{N}$ with
 $\vc{x}_i\in\R^{n}$; a query $\vc{q}\in\R^{n}$; a metric $d$; an integer
 $k\ge1$. \quad\emph{Output:} a predicted label $\hat{y}$.}
  \item \kw{for} $i=1$ \kw{to} $N$ \kw{do} \hfill\textit{\small Step 1: measure}
  \item \ind $\delta_i \;\gets\; d(\vc{q},\vc{x}_i)$
  \item \kw{end for}
  \item let $\pi$ be a permutation of $\{1,\dots,N\}$ with
        $\delta_{\pi(1)}\le\delta_{\pi(2)}\le\cdots\le\delta_{\pi(N)}$
        \hfill\textit{\small Step 2: select}
  \item $\Nk(\vc{q}) \;\gets\; \bigl\{\pi(1),\ \pi(2),\ \dots,\ \pi(k)\bigr\}$
  \item \kw{for} each label $c$ occurring in $\D$ \kw{do}
        \hfill\textit{\small Step 3: vote}
  \item \ind $\displaystyle
        \mathrm{votes}(c) \;\gets\;
        \sum_{i\,\in\,\Nk(\vc{q})} \mathbf{1}\bigl[\,y_i = c\,\bigr]$
  \item \kw{end for}
  \item $\displaystyle \hat{y} \;\gets\; \argmax_{c}\ \mathrm{votes}(c)$
  \item \kw{return} $\hat{y}$
\end{algo}

\noindent Two lines deserve a second look.

Line~4 asks for a permutation that sorts the distances, which is more than the
algorithm needs --- only the first $k$ entries of $\pi$ are ever used. Line~4 is
written this way because it states \emph{what} is wanted with no commitment to
\emph{how}, and \S7.2 shows the choice of how is worth a factor of $\log N$.

Line~9 hides a decision. If two labels tie for the maximum, $\argmax$ does not
say which to return. \S8.4 supplies the usual conventions.

\begin{rmk}[The regression variant]
If the labels $y_i$ are real numbers rather than categories, the vote in lines
$6$--$9$ is replaced by an average:
\[
  \hat{y} \;=\; \frac{1}{k}\sum_{i\,\in\,\Nk(\vc{q})} y_i .
\]
Nothing else changes. Steps 1 and 2 --- which are the expensive ones --- are
identical.
\end{rmk}

%------------------------------------------------------------------------------
\subsection{A refinement: weighting by distance}
%------------------------------------------------------------------------------
The vote of \S6.3 treats all $k$ neighbours equally, so a point on the far edge
of the neighbourhood counts as much as one sitting almost on top of the query.
A common fix weights each vote by $1/\delta_i$, or by $1/\delta_i^{2}$:
\[
  \mathrm{votes}(c)
  \;=\; \sum_{i\,\in\,\Nk(\vc{q})}
        \frac{\mathbf{1}[\,y_i=c\,]}{\delta_i} .
\]
This costs $k$ extra divisions and nothing else, so it does not change any
complexity in \S7. It also softens the choice of $k$: distant neighbours are
admitted to the vote but do not count for much, so the answer changes less
abruptly as $k$ grows.

%------------------------------------------------------------------------------
\subsection{A worked example}
%------------------------------------------------------------------------------
\begin{ex}
Work in $\R^{2}$ with Euclidean distance. The dataset is eight points, four in
class $A$ and four in class $B$:
\[
  \begin{aligned}
    A:\quad & \vc{a}_1=[\,1,1\,],\quad \vc{a}_2=[\,2,3\,],\quad
              \vc{a}_3=[\,2,1\,],\quad \vc{a}_4=[\,4,5\,], \\
    B:\quad & \vc{b}_1=[\,6,4\,],\quad \vc{b}_2=[\,7,6\,],\quad
              \vc{b}_3=[\,6,8\,],\quad \vc{b}_4=[\,8,5\,].
  \end{aligned}
\]
The query is $\vc{q}=[\,5,5\,]$. Step~1 computes eight distances. It is cheaper
and equivalent to work with $\delta^{2}$ --- see \S7.5 --- so both are tabulated:

\begin{center}
\renewcommand{\arraystretch}{1.25}
\begin{tabular}{@{}l l r r l@{}}
\toprule
\textbf{Point} & \textbf{Label} & $\boldsymbol{\delta^{2}}$ & $\boldsymbol{\delta}$
  & \textbf{Rank} \\
\midrule
$\vc{a}_4=[4,5]$ & $A$ & $1$  & $1.000$ & 1 \\
$\vc{b}_1=[6,4]$ & $B$ & $2$  & $1.414$ & 2 \\
$\vc{b}_2=[7,6]$ & $B$ & $5$  & $2.236$ & 3 \\
$\vc{b}_4=[8,5]$ & $B$ & $9$  & $3.000$ & 4 \\
$\vc{b}_3=[6,8]$ & $B$ & $10$ & $3.162$ & 5 \\
$\vc{a}_2=[2,3]$ & $A$ & $13$ & $3.606$ & 6 \\
$\vc{a}_3=[2,1]$ & $A$ & $25$ & $5.000$ & 7 \\
$\vc{a}_1=[1,1]$ & $A$ & $32$ & $5.657$ & 8 \\
\bottomrule
\end{tabular}
\end{center}

\noindent All eight distances are distinct, so Step~2 has no ambiguity. Now vary
$k$:

\begin{center}
\renewcommand{\arraystretch}{1.3}
\begin{tabular}{@{}l l l l@{}}
\toprule
$\boldsymbol{k}$ & $\boldsymbol{\Nk(\vc{q})}$ & \textbf{Votes}
  & \textbf{Prediction} \\
\midrule
$1$ & $\vc{a}_4$
    & $A:1,\ B:0$ & $\hat{y}=A$ \\
$3$ & $\vc{a}_4,\vc{b}_1,\vc{b}_2$
    & $A:1,\ B:2$ & $\hat{y}=B$ \\
$5$ & $\vc{a}_4,\vc{b}_1,\vc{b}_2,\vc{b}_4,\vc{b}_3$
    & $A:1,\ B:4$ & $\hat{y}=B$ \\
\bottomrule
\end{tabular}
\end{center}

\noindent The answer at $k=1$ disagrees with the answer at $k=3$ and $k=5$, and
the table shows exactly why. The point $\vc{a}_4=[4,5]$ carries label $A$ but
sits deep inside the region occupied by $B$ --- an outlier, or a mislabelling,
or a genuinely unusual case. Being nearest to $\vc{q}$, it decides the vote by
itself when $k=1$. As soon as two more neighbours are admitted they are both
$B$, and they outvote it. \hfill$\blacksquare$
\end{ex}

\begin{rmk}
This example is the argument for $k>1$ in miniature. Raising $k$ does not make
the classifier smarter; it makes it \emph{less sensitive to any single training
point}. That is worth having when the data contains noise, and it is exactly
what you do not want when the data is clean and the true boundary is intricate.
\S8.1 is that trade-off, stated.
\end{rmk}

%==============================================================================
\section{What it costs}
%==============================================================================
Now the question Part 1 taught us to ask. We count the work for \emph{one}
query against a dataset of $N$ points in $\R^{n}$, in FLOPs (\S4.6).

%------------------------------------------------------------------------------
\subsection{Step 1: the distances}
%------------------------------------------------------------------------------
One Euclidean distance $d(\vc{q},\vc{x}_i)=\norm{\vc{q}-\vc{x}_i}$ is two
operations we have already priced. Forming the difference $\vc{q}-\vc{x}_i$ is
vector subtraction, $n$ FLOPs by \S1.1. Taking its norm is $2n$ FLOPs by \S1.4
--- $n$ mults to square, $n-1$ adds to total, one square root. So
\[
  \text{one distance} \;=\; n + 2n \;=\; 3n \text{ FLOPs} ,
\]
and Step~1 performs $N$ of them:
\[
  \boxed{\;\text{Step 1} \;=\; 3Nn \text{ FLOPs} \;=\; O(Nn) . \;}
\]

\begin{rmk}
The count is linear in \emph{both} inputs, and both matter. Doubling the number
of training points doubles the work; so does doubling the number of features.
This is the sense in which $k$-NN has no free lunch: every training point is
touched, every time, at full length.
\end{rmk}

%------------------------------------------------------------------------------
\subsection{Step 2: the selection}
%------------------------------------------------------------------------------
Step~2 must find the $k$ smallest of $N$ numbers. There are three standard ways,
and they are not equally good.

\begin{center}
\renewcommand{\arraystretch}{1.35}
\begin{tabular}{@{}l l l@{}}
\toprule
\textbf{Method} & \textbf{Cost} & \textbf{Comment} \\
\midrule
Sort all $N$ distances & $O(N\log N)$
  & what line 4 literally says; wasteful \\
Keep a heap of the $k$ best & $O(N\log k)$
  & one pass, $k$ items retained \\
Quickselect & $O(N)$ average
  & $O(N^{2})$ worst case \\
\bottomrule
\end{tabular}
\end{center}

\noindent The middle row is what libraries do. Scan the $N$ distances once,
maintaining a structure holding the best $k$ seen so far; each new distance
costs at most $O(\log k)$ to insert or reject. Since $k$ is typically a small
constant --- $3$, $5$, $11$ --- and $N$ may be in the millions, the difference
between $\log k$ and $\log N$ is the difference between a factor of $2$ and a
factor of $20$.

%------------------------------------------------------------------------------
\subsection{Step 3: the vote}
%------------------------------------------------------------------------------
Counting the labels of $k$ points and taking the largest count is $O(k)$, and
$k\le N$. Beside Steps 1 and 2 this is free, and it will not appear again.

%------------------------------------------------------------------------------
\subsection{The total}
%------------------------------------------------------------------------------
Adding the three steps,
\[
  T(N,n,k) \;=\; \underbrace{O(Nn)}_{\text{distances}}
              +\ \underbrace{O(N\log k)}_{\text{selection}}
              +\ \underbrace{O(k)}_{\text{vote}}
  \;=\; \boxed{\;O\bigl(Nn + N\log k\bigr)\;}
\]
per query. In almost every real setting the first term dominates: $n$ is the
number of features, often hundreds or thousands, while $\log k$ is a small
number like $2$ or $3$. So the working summary is that a $k$-NN query costs
$O(Nn)$, and $M$ queries cost $O(MNn)$.

\begin{center}
\renewcommand{\arraystretch}{1.35}
\begin{tabular}{@{}l l l@{}}
\toprule
\textbf{Quantity} & \textbf{Cost} & \textbf{Why} \\
\midrule
Training                  & $O(1)$      & store the data; nothing is computed \\
Memory                    & $O(Nn)$     & every training point must be kept \\
One query                 & $O(Nn + N\log k)$ & \S7.4 \\
$M$ queries               & $O(MNn)$    & no work is shared between queries \\
\bottomrule
\end{tabular}
\end{center}

\begin{rmk}[A lazy learner]
The first row is the strangest entry in the table and the most characteristic.
Training a $k$-NN classifier costs nothing: you keep the data. Algorithms that
behave this way are called \textbf{lazy} --- they defer all computation to query
time --- as opposed to \textbf{eager} learners, which do expensive work once up
front and then answer queries cheaply.

Neither is better in the abstract, and the choice is a choice about where you
can afford to spend. Laziness is excellent when the dataset changes constantly,
because there is no model to rebuild; adding a training point is free. It is
terrible when queries are frequent and the dataset is large, because you pay
$O(Nn)$ every single time and never amortise it. \S10.2 is about buying back
some of that cost with a one-off index.
\end{rmk}

%------------------------------------------------------------------------------
\subsection{One saving, and it is Part 1's saving}
%------------------------------------------------------------------------------
There is a cheap improvement to Step~1 that costs nothing to adopt. The
algorithm never uses the \emph{value} of a distance --- only the \emph{order} of
the distances. And squaring is increasing on non-negative numbers, so for
$\delta,\delta'\ge0$,
\[
  \delta \le \delta' \iff \delta^{2} \le \delta'^{2} .
\]
The ranking by $\delta$ and the ranking by $\delta^{2}$ are therefore the same
ranking, and we may compare squared distances and never take a square root:
\[
  \delta_i^{2} = \sum_{j=1}^{n}(q_j - x_{ij})^{2}
  \qquad\text{costs } n + n + (n-1) = 3n-1 \text{ FLOPs} .
\]
The saving is one square root per training point, so $N$ square roots per query.

\begin{rmk}
This is exactly the second saving of \S1.6 --- ``do not take a square root you
are going to square'' --- arriving in a new setting, and it is worth noticing
that the argument is not the same one. In \S1.6 the square root cancelled
algebraically. Here it does not cancel at all; it is discarded because the
answer only depends on the ordering, and the square root preserves orderings.
Different reason, same square root, same saving.

In Big-O terms nothing has happened: $3n-1$ and $3n$ are both $O(n)$. In wall
clock terms something has, because a square root is among the slowest operations
a processor offers (\S1.2), and this removes one per training point per query.
It is the standard implementation, and the table in \S6.5 was built on it.
\end{rmk}

%==============================================================================
\section{Choosing \texorpdfstring{$k$}{k}, and what goes wrong}
%==============================================================================
%------------------------------------------------------------------------------
\subsection{Small $k$, large $k$}
%------------------------------------------------------------------------------
The example of \S6.5 showed the two failure directions in one dataset.

\begin{center}
\renewcommand{\arraystretch}{1.35}
\begin{tabular}{@{}l l l@{}}
\toprule
& \textbf{$k$ too small} & \textbf{$k$ too large} \\
\midrule
Boundary shape   & jagged, follows every point & smooth, ignores real detail \\
Sensitive to     & noise and mislabelling      & nothing much \\
Failure mode     & memorises the training data & predicts the majority label \\
Standard name    & \emph{overfitting}          & \emph{underfitting} \\
\bottomrule
\end{tabular}
\end{center}

\noindent There is no formula for the right $k$. It is chosen empirically: hold
back part of the labelled data, run the classifier on it for several values of
$k$, and keep the value that gets the most answers right. Doing this properly
--- so that the held-back data is not quietly used to make the choice it is
supposed to be judging --- is called \emph{cross-validation}, and it belongs to
a statistics course rather than this one. A common starting guess is
$k\approx\sqrt{N}$.

%------------------------------------------------------------------------------
\subsection{Feature scaling, and why it is not optional}
%------------------------------------------------------------------------------
This is the failure that catches everyone once. Suppose houses are described by
\[
  \vc{x} = [\ \underbrace{\text{area in m}^{2}}_{\text{roughly }50\text{--}300},
             \ \underbrace{\text{number of bedrooms}}_{\text{roughly }1\text{--}5}\ ] .
\]
Take two houses of equal area differing by $4$ bedrooms, and two houses with the
same bedrooms differing by $4$ square metres. Squared Euclidean distance scores
both differences as $16$. But one pair is a studio next to a family home and the
other pair is indistinguishable.

The arithmetic is doing precisely what it was told. Euclidean distance adds the
squared differences of the components with equal weight, so a feature measured
in large units contributes large differences and dominates the sum. Change the
area to square centimetres and it will dominate absolutely; change it to
hectares and it will vanish. \emph{The units of your features silently set their
importance.}

\begin{defn}[Standardization]
To \textbf{standardize} a dataset, replace each feature by
\[
  x_{ij} \;\longmapsto\; \frac{x_{ij}-\mu_j}{\sigma_j} ,
\]
where $\mu_j$ is the mean of feature $j$ across all $N$ points and $\sigma_j$ is
its standard deviation. Every feature then has mean $0$ and spread $1$, and none
is privileged by its units.
\end{defn}

\begin{warn}[Distance is not unit-free]
Nothing in the metric axioms of the vector unit, \S14.1, protects you here.
$d_2$ is a perfectly good metric on $\R^{n}$ whatever the columns mean; the
axioms constrain how $d$ behaves, not whether the numbers it is fed are
comparable. Making them comparable is your job, it happens before the algorithm
runs, and skipping it does not produce an error --- it produces a classifier
that quietly ignores most of your features.
\end{warn}

\begin{rmk}
An alternative to rescaling is to change the metric. Instead of weighting every
feature equally, use
$d(\vc{x},\vc{y})^{2}=\sum_j w_j (x_j-y_j)^{2}$ with weights $w_j$ chosen to
reflect how much each feature should matter. Standardization is the special case
$w_j = 1/\sigma_j^{2}$, which says ``every feature matters equally once its own
spread is accounted for''.
\end{rmk}

%------------------------------------------------------------------------------
\subsection{The curse of dimensionality}
%------------------------------------------------------------------------------
The third difficulty is subtler and it is specific to large $n$. As the
dimension grows, the distances from a query to all the training points tend to
become nearly equal --- the nearest point is barely nearer than the farthest.
When that happens the phrase ``the $k$ nearest neighbours'' stops picking out
anything meaningful: the ranking still exists, but it is decided by noise.

The intuition is a volume count. In $\R^{n}$ a ball of radius $r$ has volume
proportional to $r^{n}$, so a small increase in radius sweeps up a huge
additional region. Points spread through a high-dimensional space are
overwhelmingly likely to sit at similar middling distances from each other, and
close neighbours become vanishingly rare unless $N$ grows exponentially in $n$.

\begin{rmk}
This is not an argument against $k$-NN, which remains one of the most useful
methods there is on high-dimensional data --- image and text retrieval are
exactly that. It is an argument for not using the \emph{raw} dimension. The
standard response is to reduce the dimension first, mapping the data into a
space of far fewer coordinates that preserves the distances that matter, and
then run $k$-NN there. The tool for doing that linearly is the singular value
decomposition, which is a later unit of this course.
\end{rmk}

%------------------------------------------------------------------------------
\subsection{Ties}
%------------------------------------------------------------------------------
Two kinds of tie can occur, and both need a rule stated in advance.

\begin{itemize}[leftmargin=1.4em,itemsep=4pt]
  \item \textbf{A tie in distance} at the boundary of the neighbourhood: more
        than $k$ points are within the cutoff. Conventions: take all of them;
        take the ones with the lowest index; or break the tie at random. Any is
        defensible, and libraries differ, so it must be checked rather than
        assumed.
  \item \textbf{A tie in the vote}: two labels receive the same count. The usual
        rules are to fall back to the label of the single nearest neighbour, to
        decrease $k$ by one and re-vote, or --- best --- to use distance
        weighting (\S6.4), under which exact ties are essentially impossible.
        Choosing $k$ odd removes the problem for two classes and does nothing
        for three.
\end{itemize}

\begin{rmk}
It is worth remembering, from Part 1, that these are floating-point distances
(\S3). Two distances that are equal in exact arithmetic may differ in the last
bit after rounding, and two that differ mathematically may round to the same
float. So a tie-breaking rule is not only a modelling convention; it is also
what keeps the output from depending on rounding.
\end{rmk}

%==============================================================================
\section{\texorpdfstring{$k$}{k}-NN and clustering}
%==============================================================================
%------------------------------------------------------------------------------
\subsection{Two different problems}
%------------------------------------------------------------------------------
\begin{warn}[$k$-NN does not cluster]
The phrase ``$k$-NN clustering'' is in wide circulation and it conflates two
genuinely different problems.

\emph{$k$-NN} is \textbf{supervised}. It requires a labelled dataset, and it
answers a question about one new point: what label should $\vc{q}$ get? It never
partitions anything.

\emph{Clustering} is \textbf{unsupervised}. There are no labels. The task is to
partition the dataset itself into groups of mutually similar points, and the
groups are the output. That is Part 3's subject, $k$-means.

The two are neighbours in every syllabus because both run on nothing but
distances between feature vectors, and the $k$ in each is an integer the user
must choose --- but they are different $k$'s answering different questions.
\end{warn}

\noindent There is nevertheless a real construction that the phrase legitimately
names, and it is worth having, because it is the substrate of several genuine
clustering algorithms.

%------------------------------------------------------------------------------
\subsection{The $k$-NN graph}
%------------------------------------------------------------------------------
\begin{defn}[$k$-nearest-neighbour graph]
Let $P=\{\vc{p}_1,\dots,\vc{p}_N\}\subseteq\R^{n}$ and fix a metric $d$ and an
integer $k$. The \textbf{$k$-NN graph} of $P$ has the points of $P$ as its
vertices, and joins $\vc{p}_i$ to $\vc{p}_j$ whenever $\vc{p}_j$ is one of the
$k$ nearest neighbours of $\vc{p}_i$, or $\vc{p}_i$ is one of the $k$ nearest
neighbours of $\vc{p}_j$.
\end{defn}

\noindent The construction is $k$-NN run $N$ times, once with each point of $P$
as the query and the rest as the dataset --- so it costs $O(N^{2}n)$ by \S7.4,
with no labels involved anywhere. What comes out is not a classification but a
\emph{structure}: which points are near which.

Now cluster the graph rather than the points. The simplest rule is to take the
\textbf{connected components} --- the groups of vertices reachable from one
another by following edges --- and call each one a cluster. Two points end up
together when a chain of short hops connects them, even if they are far apart
directly, which is the essential difference from $k$-means and is why this
family can find long, curved clusters that $k$-means cannot.

%------------------------------------------------------------------------------
\subsection{A worked example}
%------------------------------------------------------------------------------
\begin{ex}
Six unlabelled points in $\R^{2}$:
\[
  \vc{p}_1=[\,0,0\,],\quad \vc{p}_2=[\,1,0\,],\quad \vc{p}_3=[\,0,2\,],
  \quad
  \vc{p}_4=[\,10,10\,],\quad \vc{p}_5=[\,11,10\,],\quad \vc{p}_6=[\,10,12\,].
\]
Build the $1$-NN graph under Euclidean distance. For each point, find its single
nearest neighbour:

\begin{center}
\renewcommand{\arraystretch}{1.25}
\begin{tabular}{@{}l l l l@{}}
\toprule
\textbf{Point} & \textbf{Distances to the others} & \textbf{Nearest}
  & \textbf{Edge} \\
\midrule
$\vc{p}_1$ & $d(\vc{p}_2)=1$, $d(\vc{p}_3)=2$, rest $\ge 14.1$
           & $\vc{p}_2$ & $\{\vc{p}_1,\vc{p}_2\}$ \\
$\vc{p}_2$ & $d(\vc{p}_1)=1$, $d(\vc{p}_3)=\sqrt5\approx2.24$, rest $\ge 13.5$
           & $\vc{p}_1$ & $\{\vc{p}_1,\vc{p}_2\}$ \\
$\vc{p}_3$ & $d(\vc{p}_1)=2$, $d(\vc{p}_2)\approx2.24$, rest $\ge 12.8$
           & $\vc{p}_1$ & $\{\vc{p}_1,\vc{p}_3\}$ \\
$\vc{p}_4$ & $d(\vc{p}_5)=1$, $d(\vc{p}_6)=2$, rest $\ge 12.8$
           & $\vc{p}_5$ & $\{\vc{p}_4,\vc{p}_5\}$ \\
$\vc{p}_5$ & $d(\vc{p}_4)=1$, $d(\vc{p}_6)=\sqrt5\approx2.24$, rest $\ge 14.9$
           & $\vc{p}_4$ & $\{\vc{p}_4,\vc{p}_5\}$ \\
$\vc{p}_6$ & $d(\vc{p}_4)=2$, $d(\vc{p}_5)\approx2.24$, rest $\ge 15.6$
           & $\vc{p}_4$ & $\{\vc{p}_4,\vc{p}_6\}$ \\
\bottomrule
\end{tabular}
\end{center}

\noindent The distinct edges are
$\{\vc{p}_1,\vc{p}_2\}$, $\{\vc{p}_1,\vc{p}_3\}$,
$\{\vc{p}_4,\vc{p}_5\}$, $\{\vc{p}_4,\vc{p}_6\}$, and no edge joins the two
groups, because every cross distance exceeds $12$ while every within-group
distance is at most $2.24$.

\begin{center}
\begin{tikzpicture}[scale=0.40]
  \draw[axisline] (-1.2,0) -- (14,0) node[right,font=\scriptsize,black]{$x_1$};
  \draw[axisline] (0,-1.2) -- (0,14) node[above,font=\scriptsize,black]{$x_2$};
  \draw[RoyalBlue!70,very thick] (0,0) -- (1,0);
  \draw[RoyalBlue!70,very thick] (0,0) -- (0,2);
  \draw[BrickRed!70,very thick] (10,10) -- (11,10);
  \draw[BrickRed!70,very thick] (10,10) -- (10,12);
  \foreach \x/\y in {0/0, 1/0, 0/2}
     { \node[circle,fill=RoyalBlue!80,draw=RoyalBlue,inner sep=0pt,
             minimum size=6pt] at (\x,\y) {}; }
  \foreach \x/\y in {10/10, 11/10, 10/12}
     { \node[circle,fill=BrickRed!80,draw=BrickRed,inner sep=0pt,
             minimum size=6pt] at (\x,\y) {}; }
  \node[below left,font=\scriptsize]  at (0.1,0.1)  {$\vc{p}_1$};
  \node[below,font=\scriptsize]       at (1.3,-0.1) {$\vc{p}_2$};
  \node[left,font=\scriptsize]        at (-0.2,2)   {$\vc{p}_3$};
  \node[below left,font=\scriptsize]  at (10.1,10.1){$\vc{p}_4$};
  \node[right,font=\scriptsize]       at (11.3,10)  {$\vc{p}_5$};
  \node[above,font=\scriptsize]       at (10,12.5)  {$\vc{p}_6$};
  \node[font=\scriptsize,RoyalBlue] at (3.6,1.6) {component 1};
  \node[font=\scriptsize,BrickRed]  at (6.6,11.4) {component 2};
\end{tikzpicture}
\end{center}

\noindent The graph has exactly two connected components,
$\{\vc{p}_1,\vc{p}_2,\vc{p}_3\}$ and $\{\vc{p}_4,\vc{p}_5,\vc{p}_6\}$, and those
are the two clusters. No labels were used and none were available: this is
unsupervised, and it is what ``$k$-NN clustering'' properly refers to.
\hfill$\blacksquare$
\end{ex}

\begin{rmk}
Connected components are the crudest way to cut a $k$-NN graph and are shown
here because they can be read off by eye. Real methods cut it more carefully:
\emph{spectral clustering} uses the eigenvectors of a matrix built from the
graph --- a technique this course will be able to explain once eigenvalues
arrive --- and \emph{DBSCAN} and \emph{HDBSCAN} use neighbourhood density to
decide which edges to keep. \emph{UMAP}, the standard tool for visualising
high-dimensional data, begins by building exactly this graph. All of them start
where \S9.2 starts.
\end{rmk}

%==============================================================================
\section{In practice}
%==============================================================================
%------------------------------------------------------------------------------
\subsection{Where it is used}
%------------------------------------------------------------------------------
\begin{center}
\renewcommand{\arraystretch}{1.4}
\begin{tabular}{@{}l >{\raggedright\arraybackslash}p{9.4cm}@{}}
\toprule
\textbf{Application} & \textbf{The vectors, and the metric} \\
\midrule
Recommendation
  & Each user is a vector of ratings; neighbours are users with similar taste,
    usually under cosine distance, and their preferences are the prediction. \\
Image retrieval
  & An image becomes a vector via a neural network; ``find similar images'' is
    literally a nearest-neighbour query. \\
Semantic search, RAG
  & Text passages are embedded as vectors; a question is embedded the same way,
    and the $k$ nearest passages are retrieved and handed to a language model.
    This is now among the largest uses of nearest-neighbour search anywhere. \\
Anomaly detection
  & The distance to the $k$-th nearest neighbour is an outlier score: points in
    sparse regions are far from everything. \\
Handwriting recognition
  & The classic benchmark. Plain $k$-NN with Euclidean distance on raw
    $28\times28$ pixel vectors classifies handwritten digits correctly about
    $97\%$ of the time --- with no training at all, which is why it remains the
    baseline every new method is expected to beat. \\
\bottomrule
\end{tabular}
\end{center}

\noindent The pattern across the table is worth naming. $k$-NN is chosen when
similarity is the whole question, when the data changes often enough that
retraining a model would be a nuisance, and when an answer needs to be
explainable --- ``we predicted this because of those five specific cases'' is a
justification a person can check.

%------------------------------------------------------------------------------
\subsection{Making it fast}
%------------------------------------------------------------------------------
Everything in \S7 assumed the obvious implementation: compare the query against
all $N$ points. That is called \textbf{brute force} or exhaustive search, and
its $O(Nn)$ per query is the thing every optimisation attacks. The methods
divide into two families by whether they still return the true answer.

\subsubsection*{Exact methods}
These return exactly what \S6.3 would, faster, by ruling out regions of space
without examining the points in them.

\begin{itemize}[leftmargin=1.4em,itemsep=4pt]
  \item \textbf{$k$-d trees.} Recursively split the data by one coordinate at a
        time, building a binary tree. A query descends the tree and prunes any
        branch whose region is too far away to contain a neighbour. Queries cost
        about $O(\log N)$ in low dimension. The catch is that the pruning stops
        working as $n$ grows --- by roughly $n\approx20$ the tree examines
        almost everything and is slower than brute force, for the reason given
        in \S8.3.
  \item \textbf{Ball trees.} Split by enclosing balls rather than coordinate
        planes. They tolerate higher dimension better than $k$-d trees and work
        with any metric, not only Euclidean.
  \item \textbf{Triangle-inequality pruning.} If $d(\vc{q},\vc{x})$ is known and
        $d(\vc{x},\vc{y})$ was precomputed, then \textbf{M4} of the vector unit,
        \S14.1, gives $d(\vc{q},\vc{y}) \ge d(\vc{q},\vc{x}) - d(\vc{x},\vc{y})$,
        and if that lower bound already exceeds the current $k$-th best
        distance, $\vc{y}$ can be discarded without ever being measured. This is
        a metric axiom doing computational work, and Part 3 will use the same
        trick on $k$-means.
\end{itemize}

\subsubsection*{Approximate methods (ANN)}
Above a few dozen dimensions, exact methods lose to brute force, and the
practical answer is to stop insisting on exactness. \textbf{Approximate nearest
neighbour} search returns neighbours that are \emph{almost always} the true
ones, thousands of times faster.

\begin{itemize}[leftmargin=1.4em,itemsep=4pt]
  \item \textbf{LSH} (locality-sensitive hashing) hashes vectors so that nearby
        points collide with high probability; a query examines only its own
        bucket.
  \item \textbf{HNSW} (hierarchical navigable small world) builds a layered
        proximity graph --- a $k$-NN graph, \S9.2, with long-range shortcuts ---
        and answers a query by greedy descent. It is the current default in most
        systems.
  \item \textbf{IVF with product quantization} partitions the space into cells,
        searches only the nearest few, and stores each vector in compressed form
        so that millions fit in memory.
\end{itemize}

\begin{rmk}
The trade is explicit and it is measured. \emph{Recall@$k$} is the fraction of
the true $k$ neighbours that a method actually returns, and ANN systems are
tuned to a target such as $95\%$ recall at some queries-per-second budget.
Whether losing one neighbour in twenty matters depends entirely on the
application --- for retrieving passages to answer a question it is invisible;
for a legal or medical audit trail it may not be acceptable.
\end{rmk}

%------------------------------------------------------------------------------
\subsection{What the industry uses}
%------------------------------------------------------------------------------
\begin{center}
\renewcommand{\arraystretch}{1.35}
\begin{tabular}{@{}l >{\raggedright\arraybackslash}p{9.0cm}@{}}
\toprule
\textbf{Tool} & \textbf{What it is} \\
\midrule
\texttt{scikit-learn}
  & \texttt{KNeighborsClassifier}. Its \texttt{algorithm="auto"} chooses
    between brute force, a $k$-d tree and a ball tree from $N$ and $n$. The
    teaching and prototyping standard. \\
\texttt{FAISS}
  & Meta's library for large-scale similarity search, and the reference
    implementation: IVF, product quantization, GPU support. \\
\texttt{hnswlib}, \texttt{ScaNN}, \texttt{Annoy}
  & Focused single-index ANN implementations, the latter two from Google and
    Spotify. \\
Vector databases
  & Pinecone, Weaviate, Milvus, Qdrant, and the \texttt{pgvector} extension to
    PostgreSQL. These are ANN indexes wrapped in a database: storage,
    filtering, updates and a query API. \\
\bottomrule
\end{tabular}
\end{center}

\noindent Two things are worth noticing about that table. The first is that
almost all of it is recent --- the vector-database category barely existed
before large language models made embedding-and-retrieving the standard way to
give a model access to documents. The second is that every entry is an
implementation of \S6.3. The rule being computed is still: measure the
distances, take the $k$ smallest, and let them vote. Everything else is an index
built to avoid measuring all $N$ of them.

\lookahead Part 3 turns to the other problem. There will be no labels, so there
is nothing to vote; the task is to find the groups in the data rather than to
place a new point into known ones. The tools are the same --- distances between
vectors, and counting the cost of computing them --- and one new object is
needed, the mean of a set of vectors.

\end{document}
